\documentclass[11pt]{article}
\usepackage[margin=1.1in]{geometry}
\usepackage{amsmath,amssymb,amsthm,mathtools,booktabs,array,tabularx,longtable}
\usepackage[dvipsnames]{xcolor}
\usepackage[colorlinks=true,linkcolor=MidnightBlue,citecolor=MidnightBlue,urlcolor=MidnightBlue]{hyperref}
\usepackage{fancyhdr}
\pagestyle{fancy}\fancyhf{}\renewcommand{\headrulewidth}{0pt}
\fancyfoot[R]{\thepage}
\fancyfoot[L]{\scriptsize\color{gray}Generated by the Transport Lab from approved parts of its store; not a deposited paper. Papers of record: \href{https://zenodo.org/communities/transport-papers}{Zenodo, Transport Papers}}
\fancypagestyle{plain}{\fancyhf{}\fancyfoot[R]{\thepage}\fancyfoot[L]{\scriptsize\color{gray}Generated by the Transport Lab from approved parts of its store; not a deposited paper. Papers of record: \href{https://zenodo.org/communities/transport-papers}{Zenodo, Transport Papers}}\renewcommand{\headrulewidth}{0pt}}
\newtheorem{theorem}{Theorem}[section]
\newtheorem{proposition}[theorem]{Proposition}
\newtheorem{lemma}[theorem]{Lemma}
\newtheorem{corollary}[theorem]{Corollary}
\newtheorem{conjecture}[theorem]{Conjecture}
\newtheorem{refutation}[theorem]{Refutation}
\newtheorem{prediction}[theorem]{Prediction}
\newtheorem{principle}[theorem]{Principle}
\newtheorem{observation}[theorem]{Observation}
\newtheorem{axiom}[theorem]{Axiom}
\theoremstyle{definition}
\newtheorem{definition}[theorem]{Definition}
\newtheorem{fact}[theorem]{Fact}
\newtheorem{assembled}[theorem]{Assembly}
\newtheorem{postulate}[theorem]{Postulate}
\newtheorem{premise}[theorem]{Premise}
\newtheorem{hypothesis}[theorem]{Hypothesis}
\newtheorem{remark}[theorem]{Remark}
\newtheorem{example}[theorem]{Example}
\theoremstyle{plain}
\setlength{\marginparwidth}{0.8in}
\newcommand{\changedmark}[2]{\marginpar{\raggedright\scriptsize\color{#1}#2}}
\providecommand{\Aa}{}\renewcommand{\Aa}{\mathcal{A}}
\providecommand{\Ad}{}\renewcommand{\Ad}{\widetilde{\mathrm{Ad}}}
\providecommand{\Aev}{}\renewcommand{\Aev}{\mathcal{A}^{+}}
\providecommand{\Aod}{}\renewcommand{\Aod}{\mathcal{A}^{-}}
\providecommand{\B}{}\renewcommand{\B}{\mathrm{B}}
\providecommand{\C}{}\renewcommand{\C}{\mathbb{C}}
\providecommand{\CAR}{}\renewcommand{\CAR}{\mathrm{CAR}}
\providecommand{\CS}{}\renewcommand{\CS}{{\mathbb{S}^{+}}}
\providecommand{\CSm}{}\renewcommand{\CSm}{{\mathbb{S}^{-}}}
\providecommand{\Cl}{}\renewcommand{\Cl}{\mathrm{Cl}}
\providecommand{\Con}{}\renewcommand{\Con}{\mathcal{C}}
\providecommand{\Dl}{}\renewcommand{\Dl}{\Delta}
\providecommand{\EN}{}\renewcommand{\EN}{E_{N}}
\providecommand{\End}{}\renewcommand{\End}{\mathrm{End}}
\providecommand{\F}{}\renewcommand{\F}{\mathbb{F}}
\providecommand{\G}{}\renewcommand{\G}{\Gamma}
\providecommand{\HH}{}\renewcommand{\HH}{\mathbb{H}}
\providecommand{\IS}{}\renewcommand{\IS}{I\sigma_3}
\providecommand{\J}{}\renewcommand{\J}{\mathbf{J}}
\providecommand{\LH}{}\renewcommand{\LH}{L_{H}}
\providecommand{\Mat}{}\renewcommand{\Mat}{\mathrm{M}}
\providecommand{\Mo}{}\renewcommand{\Mo}{M^{\circ}}
\providecommand{\Mp}{}\renewcommand{\Mp}{M^{+}}
\providecommand{\Msun}{}\renewcommand{\Msun}{M_{\odot}}
\providecommand{\Mt}{}\renewcommand{\Mt}{\widetilde{M}}
\providecommand{\Ncone}{}\renewcommand{\Ncone}{\mathcal{N}}
\providecommand{\Neff}{}\renewcommand{\Neff}{N_{\mathrm{eff}}}
\providecommand{\PCS}{}\renewcommand{\PCS}{\mathbb{P}{\mathbb{S}^{+}}}
\providecommand{\PT}{}\renewcommand{\PT}{\mathbb{PT}}
\providecommand{\Pin}{}\renewcommand{\Pin}{\mathrm{Pin}}
\providecommand{\Q}{}\renewcommand{\Q}{\mathbb{Q}}
\providecommand{\R}{}\renewcommand{\R}{\mathbb{R}}
\providecommand{\Ric}{}\renewcommand{\Ric}{\mathrm{Ric}}
\providecommand{\SBH}{}\renewcommand{\SBH}{S_{\mathrm{BH}}}
\providecommand{\Sc}{}\renewcommand{\Sc}{\mathbb{S}}
\providecommand{\Sch}{}\renewcommand{\Sch}{\mathrm{P}}
\providecommand{\Sig}{}\renewcommand{\Sig}{\Sigma}
\providecommand{\Som}{}\renewcommand{\Som}{S_\omega}
\providecommand{\Sp}{}\renewcommand{\Sp}{\mathbb{S}}
\providecommand{\Sph}{}\renewcommand{\Sph}{\mathrm{S}}
\providecommand{\Spi}{}\renewcommand{\Spi}{S_\pi}
\providecommand{\Spin}{}\renewcommand{\Spin}{\mathrm{Spin}}
\providecommand{\Sth}{}\renewcommand{\Sth}{S_{\mathrm{th}}}
\providecommand{\T}{}\renewcommand{\T}{\mathbb{T}}
\providecommand{\TH}{}\renewcommand{\TH}{T_{\mathrm{H}}}
\providecommand{\Tr}{}\renewcommand{\Tr}{\operatorname{Tr}}
\providecommand{\UHF}{}\renewcommand{\UHF}{\mathrm{UHF}}
\providecommand{\Z}{}\renewcommand{\Z}{\mathbb{Z}}
\providecommand{\Zc}{}\renewcommand{\Zc}{\mathrm{Z}}
\providecommand{\cmg}{}\renewcommand{\cmg}{\mathrm{cm^{2}\,g^{-1}}}
\providecommand{\dS}{}\renewcommand{\dS}{\mathrm{dS}}
\providecommand{\dsc}{}\renewcommand{\dsc}{\mathrm{dsc}}
\providecommand{\dtPG}{}\renewcommand{\dtPG}{dt}
\providecommand{\dtS}{}\renewcommand{\dtS}{dt_{\mathrm{S}}}
\providecommand{\ellv}{}\renewcommand{\ellv}{\ell}
\providecommand{\er}{}\renewcommand{\er}{e_r}
\providecommand{\et}{}\renewcommand{\et}{e_t}
\providecommand{\ev}{}\renewcommand{\ev}{\mathrm{ev}}
\providecommand{\gl}{}\renewcommand{\gl}{\mathfrak{gl}}
\providecommand{\hatV}{}\renewcommand{\hatV}{\widehat{V}}
\providecommand{\hb}{}\renewcommand{\hb}{\bar{h}}
\providecommand{\hf}{}\renewcommand{\hf}{\hat{f}}
\providecommand{\id}{}\renewcommand{\id}{\mathrm{id}}
\providecommand{\kpc}{}\renewcommand{\kpc}{\mathrm{kpc}}
\providecommand{\lP}{}\renewcommand{\lP}{\ell_{P}}
\providecommand{\mD}{}\renewcommand{\mD}{\mathbf{m}_{D}}
\providecommand{\mb}{}\renewcommand{\mb}{\bar m}
\providecommand{\mnu}{}\renewcommand{\mnu}{\Sigma m_{\nu}}
\providecommand{\nInf}{}\renewcommand{\nInf}{n_\infty}
\providecommand{\nO}{}\renewcommand{\nO}{n_0}
\providecommand{\ninf}{}\renewcommand{\ninf}{n_{\infty}}
\providecommand{\nz}{}\renewcommand{\nz}{n_{0}}
\providecommand{\proj}{}\renewcommand{\proj}{\operatorname{proj}}
\providecommand{\pv}{}\renewcommand{\pv}{\pi}
\providecommand{\rad}{}\renewcommand{\rad}{\mathrm{rad}}
\providecommand{\sgn}{}\renewcommand{\sgn}{\operatorname{sgn}}
\providecommand{\shA}{}\renewcommand{\shA}{\mathcal{A}}
\providecommand{\shC}{}\renewcommand{\shC}{\mathcal{C}}
\providecommand{\shE}{}\renewcommand{\shE}{\mathcal{E}}
\providecommand{\shJ}{}\renewcommand{\shJ}{\mathcal{J}}
\providecommand{\shK}{}\renewcommand{\shK}{\mathcal{K}}
\providecommand{\shL}{}\renewcommand{\shL}{L}
\providecommand{\shW}{}\renewcommand{\shW}{\mathcal{W}}
\providecommand{\sll}{}\renewcommand{\sll}{\mathfrak{sl}}
\providecommand{\so}{}\renewcommand{\so}{\mathfrak{so}}
\providecommand{\su}{}\renewcommand{\su}{\mathfrak{su}}
\providecommand{\supp}{}\renewcommand{\supp}{\operatorname{supp}}
\providecommand{\tr}{}\renewcommand{\tr}{\operatorname{tr}}
\providecommand{\uu}{}\renewcommand{\uu}{\mathfrak{u}}
\providecommand{\br}{}\renewcommand{\br}[1]{\langle #1\rangle}
\providecommand{\bra}{}\renewcommand{\bra}[1]{\langle #1|}
\providecommand{\file}{}\renewcommand{\file}[1]{\texttt{#1}}
\providecommand{\hp}{}\renewcommand{\hp}[2]{( #1 \,|\, #2 )}
\providecommand{\ip}{}\renewcommand{\ip}[2]{\langle #1, #2\rangle}
\providecommand{\ket}{}\renewcommand{\ket}[1]{|#1\rangle}
\providecommand{\rel}{}\renewcommand{\rel}[2]{D\bigl(#1 \,\|\, #2\bigr)}
\providecommand{\rev}{}\renewcommand{\rev}[1]{\widetilde{#1}}
\providecommand{\slashed}{}\renewcommand{\slashed}[1]{{#1}\mkern-10mu/}
\title{\bfseries Inside the black hole}
\author{Leon Butler\\ \normalsize\href{mailto:lbutler.math@gmail.com}{\texttt{lbutler.math@gmail.com}}\\ \normalsize ORCID: \href{https://orcid.org/0009-0006-4449-5617}{0009-0006-4449-5617}}
\date{Generated by the Transport Lab, 2 November 2026}
\begin{document}
\maketitle
\begin{abstract}
\normalsize
This trail goes inside a black hole, where the fall towards the centre becomes a passage of time. Near the centre of a non-rotating black hole the geometry is of Kasner form: one direction stretches while the other two shrink, each as a fixed power of proper time.

After rescaling by the natural clock, the rotation-gauge field settles to three pure boosts with the Kasner weights, which generate exactly the Lorentz algebra $\mathfrak{so}(1,3)$. A cosmological constant leaves that limit unchanged.

For a rotating black hole the series' rule, that trouble arises where a quantity cannot be inverted, picks out exactly the singular ring. The walk ends at the unstable inner horizon and at an open conjecture about the ring.
\end{abstract}
\newpage
\noindent \textbf{Authorship and disclosure. }The founding idea of the series is the author's: that division by zero in a Clifford algebra should be read not as a prohibition but as a transport into a higher-dimensional algebra, and that the transported division need not be accompanied by an inverse operation. The direction of investigation and the critical review of the results are likewise the author's. The mathematical development and the drafting of the text were carried out in extended collaboration with Claude Fable 5 (Anthropic), an AI system.

\noindent\fbox{\parbox{\dimexpr\linewidth-2\fboxsep-2\fboxrule}{\textbf{This is not a deposited paper.} It was generated by the Transport Lab on 2 November 2026 and assembled from approved parts of its store. It follows results across several papers of the series; each is stated with its proof and credited to the paper that first proved it. The official versions of those papers are on Zenodo, in the Transport Papers community, \url{https://zenodo.org/communities/transport-papers}.}}
\newpage

\section*{Introduction}
The walk starts with the non-rotating black hole, building from its Kasner form to the limit at the singularity, and then turns to the rotating black hole.

\section{Inside the black hole}

Each step looks a little closer at the region inside the horizon, first without rotation and then with it.

We begin with the shape of the geometry inside the horizon, where the radius $r$ runs like a clock.

In proper time $\tau$, with the singularity at $\tau = 0$, the derived metric is synchronous and diagonal, and to leading order its scale factors obey \[ \sqrt{A} \propto \tau\^{}{-1/3}, \qquad r \propto \tau\^{}{2/3} \] These are the Kasner exponents $(-\tfrac13, \tfrac23, \tfrac23)$: one direction stretches while two shrink.

\begin{lemma}[Kasner form of the interior]\label{c:interior.kasner-form-interior}
Inside a non-rotating black hole the roles of space and time swap, and the fall towards the centre becomes a passage of time. In the right time variable the gravitational field near the centre looks like a Kasner universe: one direction stretches while the other two shrink, each, to leading order, as a fixed power of time. This result gives those powers, minus one third and two thirds twice, together with the first corrections.
\emph{First proved in} T6, Lem 5.1 (\href{https://doi.org/10.5281/zenodo.22721990}{doi:10.5281/zenodo.22721990}).
\end{lemma}
\noindent\emph{Rests on:} \ref{c:interior.imported-schwarzschild-curvature-spectrum}.

An observer falling towards the centre can carry a frame of reference directions along.

The adapted frame is exactly parallel along the free-fall worldlines, at every $\tau > 0$, so carrying it into the singularity is not an asymptotic question. The expansion rates reach the Kasner exponents at the single rate $O(\tau^{2/3})$, without oscillation.

\begin{lemma}[Exact parallelism, and the convergence rate]\label{c:interior.exact-parallelism-convergence-rate}
Inside a non-rotating black hole, an observer falling towards the centre can carry a set of reference directions along. This result shows that a natural choice of such directions is carried along perfectly, with no drift at all, right up to the singularity. It also gives the exact rate at which the stretching and squeezing of space approach the idealised Kasner pattern, and shows that the approach is smooth, without oscillation.
\emph{First proved in} T6, Lem 5.2 (\href{https://doi.org/10.5281/zenodo.22721990}{doi:10.5281/zenodo.22721990}).
\end{lemma}
\noindent\emph{Rests on:} \ref{c:interior.kasner-form-interior}.

With a parallel frame in hand, the rotation-gauge field can be written down exactly.

On each unit tangential direction the field $\Omega$ is a pure boost $B_i = e_0 \wedge e_i$ at the expansion rate $\sigma_i$, plus one rotation from the round spheres: \[ \Omega(e\_1) = -\sigma\_1 B\_1, \qquad \Omega(e\_3) = -\sigma\_3 B\_3 + \frac{\cot\theta}{r}\,R\_{23} \] Along the direction of fall it vanishes, $\Omega(e_0) = 0$.

\begin{lemma}[The interior connection, exactly]\label{c:interior.interior-connection-exactly}
This result writes down exactly how directions twist into one another as one moves about inside a non-rotating black hole. Along each direction the effect is a pure stretching, except for one extra rotation term that reflects the curvature of the spheres of fixed radius. That rotation term grows like one over the radius.
\emph{First proved in} T6, Lem 5.3 (\href{https://doi.org/10.5281/zenodo.22721990}{doi:10.5281/zenodo.22721990}).
\end{lemma}
\noindent\emph{Rests on:} \ref{c:interior.exact-parallelism-convergence-rate}.

This is the central result about the approach to the singularity.

After Hubble normalisation, rescaling by the natural clock, the connection tends to \[ -\sum\_i p\_i\, B\_i \otimes \hat\theta\^{}i \] three pure boosts weighted by the Kasner exponents $p_i$, for Schwarzschild and every admissible metric. The scale tractor extends as a parallel section, and the boundary holonomy algebra is exactly $\Lambda^2 W \cong \mathfrak{so}(1,3)$.

\begin{theorem}[Weighted limit theorem]\label{c:interior.weighted-limit-theorem}
This is the central result about the approach to the singularity of a non-rotating black hole, and of a wider class of similar singularities. After rescaling by the natural clock, the rules for carrying directions settle to a clean limit at the singularity: three pure stretchings with the Kasner weights, and nothing else. These generate exactly the algebra of Lorentz transformations, while the parts of the larger symmetry that caused obstructions elsewhere in the series receive nothing.
\emph{First proved in} T6, Thm 5.6 (\href{https://doi.org/10.5281/zenodo.22721990}{doi:10.5281/zenodo.22721990}).
\end{theorem}
\noindent\emph{Rests on:} \ref{c:interior.boosts-generate-exactly-levi-factor}, \ref{c:interior.covariant-form-vacuum-equations}, \ref{c:interior.exact-parallelism-convergence-rate}, \ref{c:interior.interior-connection-exactly}, \ref{c:interior.kasner-form-interior}, \ref{c:interior.metric-parallelism-trace-schouten-condition}, \ref{c:interior.non-null-stabilisers-section-they-act}, \ref{c:interior.scale-tractor-leaves-cone}, \ref{c:interior.stabilisers-flat-model}.

A natural worry is whether that limit survives a cosmological constant.

It does: near the centre the corrections are of order $\Lambda\tau^2$, below the $O(\tau^{2/3})$ recovery rate, so the exponents, data and rate are unchanged. For $\Lambda > 0$ the limit algebra lies in the reductive stabiliser $\mathfrak{so}(1,4)$, a simplification rather than a complication.

\begin{theorem}[Schwarzschild-de Sitter weighted limit]\label{c:interior.schwarzschild-de-sitter-weighted-limit}
A cosmological constant, the ingredient behind the accelerating expansion of the universe, changes a black hole's gravitational field only very slightly near its centre. This result shows that the limiting behaviour at the singularity is the same as without it, with the same powers and the same rates. The one change is algebraic, and it makes the picture simpler rather than more complicated.
\emph{First proved in} T6, Thm 5.9 (\href{https://doi.org/10.5281/zenodo.22721990}{doi:10.5281/zenodo.22721990}).
\end{theorem}
\noindent\emph{Rests on:} \ref{c:interior.boosts-generate-exactly-levi-factor}, \ref{c:interior.exact-parallelism-convergence-rate}, \ref{c:interior.kasner-form-interior}, \ref{c:interior.non-null-stabilisers-section-they-act}, \ref{c:interior.scale-tractor-leaves-cone}, \ref{c:interior.weighted-limit-theorem}.

The walk now turns to rotating black holes.

In the commuting subalgebra $\R \oplus \R I$, with $I$ the pseudoscalar and $I^2 = -1$, the Kerr curvature is the Schwarzschild shape operator times the scale \[ $\mathfrak{M}$ = -M\,(r - Ia\cos\theta)\^{}{-3} \] Its base fails to be invertible exactly on the ring $r = 0$, $\theta = \pi/2$, which is where the black hole is singular.

\begin{proposition}[Complex factorisation and the location of the ring]\label{c:interior.complex-factorisation-location-ring}
The curvature of a rotating black hole can be written as the curvature of a non-rotating one, with the radius replaced by a complex number built from the radius and the angle. The series' rule is that trouble arises where a quantity cannot be inverted. Applied to this complex number, the rule picks out exactly the ring where a rotating black hole is singular, and says that the disc spanning the ring is regular.
\emph{First proved in} T6, Prop 6.4 (\href{https://doi.org/10.5281/zenodo.22721990}{doi:10.5281/zenodo.22721990}).
\end{proposition}
\noindent\emph{Rests on:} \ref{c:interior.exact-scale-decoupling}, \ref{c:interior.imported-kerr-curvature}.

That scalar-plus-pseudoscalar scale turns out to be a well-known object.

It is the inverse cube of the Walker-Penrose Killing-spinor scalar, \[ \mathfrak M = -M\,\chi\^{}{-3}, \qquad \chi = r - Ia\cos\theta \] So going once around the ring, the curvature's duality phase turns $-3$ times as far as the spinor's.

\begin{proposition}[The complex scale is the inverse cube of the Killing spinor]\label{c:interior.complex-scale-inverse-cube-killing-spinor}
A previous result produced a complex number that controls the curvature of a rotating black hole. This result shows that the number is, up to a cube and a constant, a well-known object: the Killing spinor behind Carter's hidden conserved quantity for orbits around a rotating black hole. Going once around the singular ring, the curvature's phase therefore turns three times as often as the spinor's, in the opposite sense.
\emph{First proved in} T6, Prop 6.8 (\href{https://doi.org/10.5281/zenodo.22721990}{doi:10.5281/zenodo.22721990}).
\end{proposition}
\noindent\emph{Rests on:} \ref{c:interior.complex-factorisation-location-ring}, \ref{c:interior.imported-kerr-curvature}, \ref{c:interior.imported-walker-penrose-killing-spinor}.

A rotating black hole also has an inner horizon, which established results show to be unstable: infalling radiation makes the curvature grow without limit there while the geometry stays continuous.

Read in the series' language, this is a type-N degeneration, the same class as at the centre of a non-rotating black hole, realised on a null hypersurface. Survival is ordered by derivative level: point-level structure extends, connection-level structure generically fails, proved in spherical symmetry and expected for Kerr.

\begin{proposition}[The filtration at the Cauchy horizon, read from the literature]\label{c:interior.filtration-cauchy-horizon-read-literature}
The inner horizon of a rotating black hole is known to be unstable: perturbations make the strength of the gravitational field blow up there, although the field itself stays continuous. This result translates those established findings into the series' language. It finds the same kind of breakdown as at the centre of a non-rotating black hole, one that generically hits the rules for comparing directions but not the field itself (proved in the spherically symmetric case and expected for rotating holes), and it places the divergence where the framework expects it.
\emph{First proved in} T6, Prop 6.11 (\href{https://doi.org/10.5281/zenodo.22721990}{doi:10.5281/zenodo.22721990}).
\end{proposition}
\noindent\emph{Rests on:} \ref{c:interior.imported-mass-inflation-stability}, \ref{c:interior.stabilisers-flat-model}, \ref{c:singularity.type-n-only-type-n}.

The walk ends with an open problem rather than a result.

Since the curvature's phase winds around the ring, the conjecture is that the completion's exceptional structure over the ring is a duality-monodromy fibration, whose structure group contains the duality circle. The winding number three of the Kerr scale would be an invariant of that fibration; it is posed, not proved.

\begin{conjecture}[Duality-monodromy fibration over the ring]\label{c:interior.duality-monodromy-fibration-over-ring}
Going once around the ring singularity of a rotating black hole, the complex phase of the curvature turns three times. The conjecture is that this twisting is the visible trace of a larger geometric structure attached to the ring, a bundle whose twisting links the different sheets of the black hole's extended spacetime. It is posed as an open problem, not proved.
\emph{First proved in} T6, Conj 6.6 (\href{https://doi.org/10.5281/zenodo.22721990}{doi:10.5281/zenodo.22721990}).
\end{conjecture}

\section{Results used from elsewhere}
The proofs above use the following results of the series, stated briefly; each is proved in the paper named.

\begin{fact}[The boosts generate exactly the Levi factor]\label{c:interior.boosts-generate-exactly-levi-factor}
The three boosts generate exactly the six-dimensional Lorentz algebra $\Lambda^2 W \cong \mathfrak{so}(1,3)$, with no component outside it.
\emph{Source:} T6, Lem 5.4 (\href{https://doi.org/10.5281/zenodo.22721990}{doi:10.5281/zenodo.22721990}).
\end{fact}

\begin{fact}[Covariant form of the vacuum equations]\label{c:interior.covariant-form-vacuum-equations}
On the invertible region of a gravitational field on flat spacetime, the vacuum field equations hold exactly when the null scale tractor $\nInf$ of the derived metric is parallel for the tractor gauge field $\Omega_{T}$.
\emph{Source:} T6, Thm 4.5 (\href{https://doi.org/10.5281/zenodo.22721990}{doi:10.5281/zenodo.22721990}).
\end{fact}

\begin{fact}[Exact scale decoupling]\label{c:interior.exact-scale-decoupling}
The Schwarzschild curvature operator is exactly $-(M/r^3)\,S$ with a fixed, trace-free shape operator $S$ satisfying $S^2 + S = 2$, so only the scale changes on the approach to the singularity.
\emph{Source:} T6, Prop 3.3 (\href{https://doi.org/10.5281/zenodo.22721990}{doi:10.5281/zenodo.22721990}).
\end{fact}

\begin{fact}[Imported: the Kerr curvature]\label{c:interior.imported-kerr-curvature}
The Kerr field strength on physical bivectors is $\mathcal{R}(B) = \tfrac{M}{2}(r - Ia\cos\theta)^{-3}(B + 3\sigma_r B\sigma_r)$, with the inverse cube taken in $\R \oplus \R I$; imported from the literature.
\emph{Source:} C. Doran and A. Lasenby, Geometric Algebra for Physicists, Cambridge University Press (2003)., Ch.~14; B. Carter, Global structure of the Kerr family of gravitational fields, Phys. Rev. 174 (1968), 1559--1571..
\end{fact}

\begin{fact}[Imported: mass inflation and $C^0$ stability]\label{c:interior.imported-mass-inflation-stability}
At the perturbed Kerr Cauchy horizon there is mass inflation and the derived metric extends continuously, and in spherical symmetry the connection generically fails to be square-integrable while tidal deformations stay finite; imported from the literature.
\emph{Source:} E. Poisson and W. Israel, Internal structure of black holes, Phys. Rev. D 41 (1990), 1796--1809.; M. Dafermos and J. Luk, The interior of dynamical vacuum black holes I: the $C^0$-stability of the Kerr Cauchy horizon, preprint, arXiv:1710.01722 (2017).; J. Luk and S.-J. Oh, Strong cosmic censorship in spherical symmetry for two-ended asymptotically flat initial data I. The interior of the black hole region, Ann. of Math. (2) 190 (2019), 1--111.; J. Luk and S.-J. Oh, Strong cosmic censorship in spherical symmetry for two-ended asymptotically flat initial data II. The exterior of the black hole region, Ann. PDE 5 (2019), Art. 6.; A. Ori, Inner structure of a charged black hole: an exact mass-inflation solution, Phys. Rev. Lett. 67 (1991), 789--792..
\end{fact}

\begin{fact}[Imported: the Schwarzschild curvature spectrum]\label{c:interior.imported-schwarzschild-curvature-spectrum}
The curvature operator $\mathcal{R}(B)$ of the Schwarzschild field is diagonal in the static frame, with eigenvalue $2M/r^3$ on the radial plane and its dual and $-M/r^3$ on the four mixed planes.
\emph{Source:} T6, Fact 2.2 (\href{https://doi.org/10.5281/zenodo.22721990}{doi:10.5281/zenodo.22721990}).
\end{fact}

\begin{fact}[Imported: the Walker-Penrose Killing spinor]\label{c:interior.imported-walker-penrose-killing-spinor}
Type-D vacua carry a Killing spinor, unique up to scale, whose Kerr invariant is proportional to $r - ia\cos\theta$, which yields Carter's constant and satisfies $\Psi_2 \propto \chi^{-3}$; imported from the literature.
\emph{Source:} M. Walker and R. Penrose, On quadratic first integrals of the geodesic equations for type $22$ spacetimes, Commun. Math. Phys. 18 (1970), 265--274.; B. Carter, Global structure of the Kerr family of gravitational fields, Phys. Rev. 174 (1968), 1559--1571..
\end{fact}

\begin{fact}[Metric parallelism, and the trace of the Schouten condition]\label{c:interior.metric-parallelism-trace-schouten-condition}
The tractor gauge field preserves the tractor metric, and a Schouten map proportional to the identity forces the factor to be $R/24$ and the field to be Einstein, so $\Sch = 0$ exactly when $\Ric = 0$.
\emph{Source:} T6, Lem 4.2 (\href{https://doi.org/10.5281/zenodo.22721990}{doi:10.5281/zenodo.22721990}).
\end{fact}

\begin{fact}[Non-null stabilisers, and the section they act on]\label{c:interior.non-null-stabilisers-section-they-act}
A non-null vector $I$ has the reductive ten-dimensional stabiliser $\mathfrak{so}(I^\perp)$, and its hyperplane section of the null quadric is de Sitter or anti-de Sitter space according to the sign of $Q(I)$.
\emph{Source:} T6, Prop 4.4 (\href{https://doi.org/10.5281/zenodo.22721990}{doi:10.5281/zenodo.22721990}).
\end{fact}

\begin{fact}[$\Lambda \neq 0$: the scale tractor leaves the cone]\label{c:interior.scale-tractor-leaves-cone}
For a gravitational field whose derived metric is Einstein with constant $\Lambda$, the section $\nInf + \tfrac{\Lambda}{6}\nO$ is parallel, every parallel section $(1, 0, \rho)$ arises this way, and its squared norm is $\Lambda/3$.
\emph{Source:} T6, Thm 4.7 (\href{https://doi.org/10.5281/zenodo.22721990}{doi:10.5281/zenodo.22721990}).
\end{fact}

\begin{fact}[Stabilisers in the flat model]\label{c:interior.stabilisers-flat-model}
In the flat model the stabiliser of the null vector $\nInf$ is the Poincar\'e algebra, that of its ray is the eleven-dimensional maximal parabolic, and the nilradical acts as translations through null shears.
\emph{Source:} T6, Prop 4.3 (\href{https://doi.org/10.5281/zenodo.22721990}{doi:10.5281/zenodo.22721990}).
\end{fact}

\begin{fact}[Type N, and only type N]\label{c:singularity.type-n-only-type-n}
The obstruction at $r = 0$ is of type N only: type C is excluded because $\Cl(1,3) \cong \Mat_2(\HH)$ is central simple, and type P imposes no transport.
\emph{Source:} T5, Prop 3.1 (\href{https://doi.org/10.5281/zenodo.22721790}{doi:10.5281/zenodo.22721790}).
\end{fact}

\section*{Conclusion}
The thread runs from the Kasner form of the non-rotating interior, through a parallel frame and an exact connection, to the limit at the singularity, and then carries the same invertibility rule into the rotating case.

No postulate of the series enters this walk; its rotating part rests on imported results, namely the Kerr curvature, its Killing spinor and the instability of its inner horizon. The inner-horizon step inherits the parts of that literature that are still conjectural for rotating black holes.

\end{document}